% Generated from the GDPR information-sheet template DOCX, then aligned to `main.tex`.
% Changes made for alignment are highlighted in red.
\documentclass[12pt]{article}

\usepackage[a4paper,margin=2.1cm]{geometry}
\usepackage[dvipsnames]{xcolor}
\usepackage{hyperref}
\usepackage{parskip}
\usepackage{lmodern}
\usepackage[T1]{fontenc}
\usepackage{enumitem}
\renewcommand{\familydefault}{\sfdefault}

\hypersetup{
  colorlinks=true,
  linkcolor=MidnightBlue,
  urlcolor=MidnightBlue,
  citecolor=MidnightBlue
}

\newcommand{\chg}[1]{\textcolor{red}{#1}}

\begin{document}

\begin{center}
  {\Large \textbf{INFORMATION SHEET FOR PARTICIPANTS}}\\[0.5em]
  {\small Ethical Clearance Reference Number: [LRU/RGO-24/25-48951]}
\end{center}

\vspace{0.5em}
\textbf{YOU WILL BE GIVEN A COPY OF THIS INFORMATION SHEET}

\vspace{1em}
\noindent\textbf{Title of project}\\
\chg{Using geolocation data to measure participation in healthy adults - validation and exploration of proxy summary measures: Onboarding and Consent}

\vspace{1.25em}
\noindent\textbf{Invitation Paragraph}\\
\chg{I would like to invite you to participate in this research project, which forms part of my Master of Science. Before you decide whether you want to take part, it is important for you to understand why the research is being done and what your participation will involve. Please take time to read the following information carefully and discuss it with others if you wish. Ask me if there is anything that is not clear or if you would like more information.}

\vspace{1.25em}
\noindent\textbf{What is the purpose of the project?}\\
\chg{This project aims to see whether geolocation (``maps'') data can help show how people take part in daily life (``participation'').}

\begin{itemize}[leftmargin=1.2em]
  \item Where someone is compared to nearby places (a relative location, e.g. 5km away rather than global coordinates)
  \item What kind of place it is (e.g. park, shop, residential address, restaurant)
  \item How often that kind of place is visited (e.g. in one week) and how long visits last
  \item How much travel is done to get there, and the type of transport used (car, public transport, walking etc.)
\end{itemize}

\vspace{1.25em}
\noindent\textbf{Why have I been invited to take part?}\\
You are being invited to participate in this project because you have expressed interest in participating in this study.

\vspace{0.75em}
\noindent\textbf{You can take part if you meet the following inclusion criteria:}
\begin{itemize}[leftmargin=1.2em]
  \item \chg{Healthy adults aged 18--60 who own a smartphone and can meet the eligibility criteria below.}
  \item \chg{You will be in the UK for the duration of data collection.}
  \item You are able to understand the consent procedure.
\end{itemize}

\vspace{0.75em}
\noindent\textbf{You must not meet any of the following exclusion criteria:}
\begin{itemize}[leftmargin=1.2em]
  \item Cognitive impairment
  \item Acute or unstable medical condition
  \item Enrolled in another clinical or interventional study that is monitoring/changing participation
  \item Non-English speaker
  \item Know or suspect you are pregnant
\end{itemize}

\vspace{1.25em}
\noindent\textbf{What will happen if I take part?}\\
\textbf{If you choose to take part in this study, here is what will happen:}

\begin{itemize}[leftmargin=1.2em]
  \item \chg{Onboarding on Day 1 (we install the app and go through how it works).}
  \item \chg{Day-to-day app use for Days 1--10.}
  \item \chg{Diary confirmation for Days 2--11 (confirm or reject what the app suggests, and add context if needed).}
  \item \chg{A short semi-structured interview after Day 11 (online or in person).}
  \item All of these activities will happen alongside your normal daily routine. You do not need to visit specific places or follow a special schedule other than the day-by-day tasks described above.
\end{itemize}

\vspace{1em}
\noindent\textbf{Do I have to take part?}\\
\chg{Participation is completely voluntary. You should only take part if you want to, and choosing not to take part will not disadvantage you in any way.}
\chg{Once you have read this information sheet, please contact us if you have any questions that would help you decide. If you decide to take part, we will ask you to sign a consent form and we will give you a copy to keep.}

\vspace{1.25em}
\noindent\textbf{Incentives}\\
None

\vspace{1.25em}
\noindent\textbf{Possible risks of taking part}\\
Participating in this study poses no additional risk compared to your normal day-to-day living.

\vspace{1.25em}
\noindent\textbf{Possible benefits of taking part}\\
There are no intended benefits of participating in this study.

\vspace{1.25em}
\noindent\textbf{Data handling and confidentiality}\\
Kings College London is the Sponsor of this research and is responsible for looking after your information.

We will need to use information from you for this research project. The research team and authorised administrative staff will use this information to do the research or to check your records to make sure that the research is being done properly.

\chg{Your data will be processed in compliance with UK data protection laws, including the UK GDPR and the Data Protection Act 2018.}

\vspace{0.75em}
\noindent\textbf{Categories of personal data}\\
\chg{We will collect your name, age, contact details, anonymised geolocation data relative to starting location, geocoding data (place type and relative location), your diary responses, and questionnaire responses from the semi-structured interview.}

\vspace{1.25em}
\noindent\textbf{Data storage}\\
Where paper records are used, the data will be stored at restricted access and secure room at the place of study. Where electronic records are used, the data will be stored on your secure mobile device before being transferred to a King's network with appropriate access controls.

\vspace{1.25em}
\noindent\textbf{Data sharing}\\
Your data will not be shared with others outside of those conducting or supervising the research. Your data will be pseudonymised whilst in use.

\vspace{0.75em}
\noindent\textbf{International data sharing}\\
Your data will not be shared with others outside of those conducting or supervising the research. Your data will be pseudonymised whilst in use.

\vspace{0.75em}
\noindent\textbf{Data retention}\\
Data storage is in line with the College's research records retention schedule. Your data may be kept for up to 5 years depending on the publication status of the research.

\vspace{0.75em}
\noindent\textbf{Anonymity in research outputs}\\
You will remain anonymous in any research outputs.

\vspace{1em}
\noindent\textbf{Archiving / future use}\\
Provide information on data archiving or storage of data for future research use (including samples if taken), including any open-science platforms and whether the data to be archived or used for future research will be anonymised data only.

\vspace{1em}
\noindent Further information on how King's processes research data can be found at:\\
\url{https://www.kcl.ac.uk/research/support/research-ethics/kings-college-london-statement-on-use-of-personal-data-in-research}

\vspace{0.75em}
\noindent Please contact the research team if you require further information on data processing or a printed version of the webpage link above.

\vspace{1.25em}
\noindent\textbf{What if I change my mind about taking part?}\\
\chg{You can withdraw at any point of the project, without having to give a reason. Withdrawing from the project will not affect you in any way. You are able to withdraw your data up until 01/05/2026, after which withdrawal of your data will no longer be possible because the data will have been anonymised or committed to the final report.}

\vspace{0.75em}
If you choose to withdraw, we will not retain the information you have provided.

\vspace{1.25em}
\noindent\textbf{What will happen to the results of the project?}\\
The results of the project will be summarised in a thesis for a Master of Sciences (Clinical Engineering) Degree. The results may be published in research journals. To obtain a copy of any published research, please contact an individual listed in the further questions section below. The anonymised data set will not be made publicly available.

\vspace{1.25em}
\noindent\textbf{Who should I contact for further information?}\\
If you have any questions or require more information about this project, please contact me using the following contact details:\\
\href{mailto:alexander.m.skondras@kcl.ac.uk}{alexander.m.skondras@kcl.ac.uk}

\vspace{1.25em}
\noindent\textbf{What if I have further questions, or if something goes wrong?}\\
If this project has harmed you in any way, or if you wish to make a complaint about the conduct of the project, you can contact King's College London using the details below for further advice and information:\\
Adam Shortland --- adam.shortland@kcl.ac.uk

\vspace{0.5em}
\noindent Thank you for reading this information sheet and for considering taking part in this research.

\end{document}

